% Einheit 2 von 3 – Spiegelebene und Spiegelgerade bestimmen (bank/spiegelung/e2.jsonl, zusammenbau v0.1)
%% TODO Zeitmarke Einheit 2: Marken-Zeile nicht lesbar
%% TODO Zweigzeile Teil 1: was hier gelernt wird (Satz in Schülersprache aus dem Einheitstitel) – Einheit 2
\einheitenkopf[e2]{Einheit 2 von 3 · Spiegelebene und Spiegelgerade bestimmen}
\zweigzeile{Abitur GK}
\setcounter{aufgabe}{12}

%% TODO Titel: Ich-kann-Satz fehlt in der Bank (Platzhalter: Kettenname) – Nr. 13
%% TODO Anweisung: ein Satz über der ersten Teilaufgabe fehlt in der Bank – Nr. 13
\begin{aufgabe}{Spiegelebene und Spiegelgerade}
%% TODO \rechenplatz{n}: Schrittzahl des Grundfalls fehlt in der Bank (nur bei fester Schreibform, 2.3 b)
\begin{teile}
\teil Was ist gesucht? „Die Punkte $A$ und $B$ liegen symmetrisch zu einer Ebene $E$. Bestimme $E$.“ \\ \kreuz{das Bild gesucht – spiegeln} \\ \kreuz{die Spiegelfläche gesucht – rückwärts}
\teil $Q(3 | 4 | 3)$ ist das Spiegelbild von $P(1 | 0 | 3)$. Bestimme die Spiegelebene $E$. $E: $ \leerfeld
\teil $g: \vec{x} = (2 | 3 | 0) + r \cdot (2 | 1 | 2)$ schneidet die $x_1x_2$-Ebene in $S(2 | 3 | 0)$. Gib eine Gleichung der an der $x_1x_2$-Ebene gespiegelten Geraden $h$ an. $h: \vec{x} = $ \leerfeld
\teil $g$ und $h$ schneiden sich in $S(2 | 1 | 3)$ und haben die Richtungsvektoren $(2 | 2 | 1)$ und $(1 | 2 | 2)$. Gib eine Ebene an, an der $g$ auf $h$ gespiegelt wird. $E: $ \leerfeld
\teil Zwei Flächen eines Körpers liegen symmetrisch zu einer Ebene $H$; $W(1 | 4 | 1)$ ist das Spiegelbild von $C(3 | 4 | 3)$, $U(3 | 0 | 1)$ liegt auf der gemeinsamen Kante. Bestimme $H$. $H: $ \leerfeld
\teil Die Abbildung zeigt die Punkte $A$, $B$ und $P$ in der $x_1x_2$-Ebene. $g$ geht durch $A$ und $P$, $g^*$ durch $B$ und $P$; $g^*$ entsteht aus $g$ durch Spiegelung an einer Geraden $h$. Zeichne $g$, $g^*$ und $h$ ein und gib einen Term an, der aus den Ortsvektoren $\vec{a}$, $\vec{b}$, $\vec{p}$ den Ortsvektor eines weiteren Punkts von $h$ liefert. \hfill (Abitur 2024 LK)

\begin{ksys}[xmin=-1,xmax=7,ymin=-1,ymax=6]\punkt{1}{1}{P}\punkt{4}{5}{A}\punkt{3}{1}{B}\end{ksys}
\end{teile}
\end{aufgabe}

%% TODO Titel: Ich-kann-Satz fehlt in der Bank (Platzhalter: Kettenname) – Nr. 14
%% TODO Anweisung: ein Satz über der ersten Teilaufgabe fehlt in der Bank – Nr. 14
\begin{aufgabe}{Spiegelebene und Spiegelgerade – Fehler finden, Begründen}
\begin{teile}
\teil Finde den Fehler. $P(1 | 4 | 2)$ und $Q(5 | 0 | 6)$ liegen symmetrisch zu $E$. \rechnung{\vec{n} &= (1 | -1 | 1) \\ E: x_1 - x_2 + x_3 &= 1 - 4 + 2 = -1}
\teil Warum ist der Verbindungsvektor von Punkt und Spiegelpunkt ein Normalenvektor der Spiegelebene?
\end{teile}
\end{aufgabe}
